\section*{第 \ref{ch:b2:reduction} 章——自同态的约化}

\begin{solution}{exo:b2:reduction:1}
$A$：$\chi_A = X^2 - 2X - 3 = (X - 3)(X + 1)$。\omterm{def:b2:reduction:eigen}{特征向量}：
对于 $3$：$(1,1)$；对于 $-1$：$(1,-1)$。于是 $P = \begin{pmatrix} 1 & 1\\ 1
& -1\end{pmatrix}$ 给出 $P^{-1}AP = \operatorname{diag}(3, -1)$。

$B = J - I$，其中 $J$ 是元素全为 1 的矩阵。$J$ 的秩为 $1$，且
$Jv = 3v$，其中 $v = (1,1,1)$；$Jw = 0$ 在平面 $x + y + z =
0$ 上成立。于是 $B$ 的谱为 $\{2, -1\}$，\omterm{def:b2:reduction:eigen}{特征子空间}为
$\operatorname{Vect}(1,1,1)$（维数 $1$）与 $\{x + y + z = 0\}$
（维数 $2$，基为 $(1,-1,0), (1,0,-1)$）。取以这三列为列的 $P$
给出 $P^{-1}BP = \operatorname{diag}(2, -1, -1)$。
\end{solution}

\begin{solution}{exo:b2:reduction:2}
\emph{\omterm{def:b2:reduction:eigen}{特征子空间}：}$\chi_C = (X-1)^2$，唯一\omterm{def:b2:reduction:eigen}{特征值}为 $1$；
$\ker(C - I) = \ker\begin{pmatrix} 0&1\\ 0&0\end{pmatrix}$ 为
直线 $\operatorname{Vect}(e_1)$：维数 $1 < 2 = m_1$，故
不\omterm{def:b2:reduction:diag}{可对角化}（\cref{thm:b2:reduction:diagcrit}）。

\emph{\omterm{def:b2:reduction:polyu}{极小多项式}：}$\mu_C$ 整除 $(X-1)^2$，且 $C \neq I$，
故 $\mu_C = (X-1)^2$：二重根，
因此不\omterm{def:b2:reduction:diag}{可对角化}（\cref{cor:b2:reduction:minpolycrit}）。
\end{solution}

\begin{solution}{exo:b2:reduction:3}
$X^2 - 5X + 6 = (X-2)(X-3)$：根为单根且完全分裂，故 $u$ \omterm{def:b2:reduction:diag}{可
对角化}（\cref{cor:b2:reduction:minpolycrit}），且
$\operatorname{Sp}(u) \subseteq \{2, 3\}$。可能的谱为
$\{2\}$（$u = 2\,\mathrm{id}$）、$\{3\}$（$u = 3\,\mathrm{id}$）或
$\{2, 3\}$。

求幂：设 $u^k = a_k\,\mathrm{id} + b_k\,u$。在\omterm{def:b2:reduction:eigen}{特征子空间}上，
此式即 $2^k = a_k + 2b_k$ 与 $3^k = a_k + 3b_k$。解得 $b_k
= 3^k - 2^k$，$a_k = 3\cdot2^k - 2\cdot 3^k$：
\[
u^k = (3\cdot 2^k - 2\cdot 3^k)\,\mathrm{id} + (3^k - 2^k)\, u .
\]
（对上述三种谱均成立，恒等式逐\omterm{def:b2:reduction:eigen}{特征值}成立。）
\end{solution}

\begin{solution}{exo:b2:reduction:4}
$u$ \omterm{def:b2:reduction:diag}{可对角化}：$P = \prod_{\lambda}(X - \lambda)$，乘积遍及谱，
该多项式零化 $u$，完全分裂且根为单根。于是 $P(u|_F) =
P(u)|_F = 0$：限制被一个完全分裂且
根为单根的多项式零化，故
\omterm{def:b2:reduction:diag}{可对角化}（\cref{cor:b2:reduction:minpolycrit}）。
\end{solution}

\begin{solution}{exo:b2:reduction:5}
$\chi_A = X^2 - X - 1$，根为 $\varphi$ 与 $\psi$（互异）：
\omterm{def:b2:reduction:diag}{可对角化}，\omterm{def:b2:reduction:eigen}{特征向量}分别为 $(\varphi, 1)$ 与 $(\psi, 1)$。
由递推关系得 $\begin{pmatrix} F_{n+1}\\ F_n \end{pmatrix} =
A^n \begin{pmatrix}1\\ 0\end{pmatrix}$。将 $(1, 0)$ 按
\omterm{def:b2:reduction:eigen}{特征向量}分解：$(1,0) = \frac{1}{\varphi - \psi}\bigl((\varphi, 1) -
(\psi, 1)\bigr)$，其中 $\varphi - \psi = \sqrt5$。用 $A^n$ 作用
则每个特征分量乘以其\omterm{def:b2:reduction:eigen}{特征值}的 $n$ 次幂；
取第二个坐标，得
\[
F_n = \frac{\varphi^n - \psi^n}{\sqrt 5} .
\]
（验证：$n = 1$ 时 $\frac{\varphi - \psi}{\sqrt5} = 1$。）
\end{solution}

\begin{solution}{exo:b2:reduction:6}
设 $P = \prod_i (X - \mu_i)$ 零化 $u^2$，完全分裂且根为单根，
其根 $\mu_i$ 即为 $u^2$ 的谱。由于 $u$ 可逆，$0$
不是 $u^2$ 的\omterm{def:b2:reduction:eigen}{特征值}（$\det u^2 = (\det u)^2 \neq 0$），故
所有 $\mu_i \neq 0$。于是
\[
Q(X) = \prod_i (X^2 - \mu_i) = \prod_i (X - \sqrt{\mu_i})(X +
\sqrt{\mu_i})
\]
零化 $u$：$\;Q(u) = \prod_i (u^2 - \mu_i\,\mathrm{id}) =
P(u^2) = 0$。其根 $\pm
\sqrt{\mu_i}$（复平方根）两两互异，因为
$\mu_i$ 互异且非零（若 $\sqrt{\mu_i} = -\sqrt{\mu_j}$
则 $\mu_i = \mu_j$）。其完全分裂且根为单根，故
$u$ \omterm{def:b2:reduction:diag}{可对角化}。

不可逆时的反例：$u = \begin{pmatrix} 0 & 1\\
0 & 0\end{pmatrix}$：$u^2 = 0$ \omterm{def:b2:reduction:diag}{可对角化}，而 $u$ 不\omterm{def:b2:reduction:diag}{可对角化}。
\end{solution}

\begin{solution}{exo:b2:reduction:7}
$A = 2I + N$，其中 $N$ 为移位矩阵（$N e_2 = e_1$，$Ne_3 = e_2$），
$N^3 = 0$，$N^2 = E_{13}$：这\emph{正是}\omterm{thm:b2:reduction:dunford}{邓福德分解}
（$2I$ 为对角矩阵，$N$ 幂零，二者可交换；唯一性使其
就是那个分解）。由于两项可交换，用二项式定理展开：
\[
A^k = 2^k I + k 2^{k-1} N + \binom k2 2^{k-2} N^2
= \begin{pmatrix}
2^k & k2^{k-1} & \binom k2 2^{k-2}\\
0 & 2^k & k2^{k-1}\\
0 & 0 & 2^k
\end{pmatrix},
\]
\[
\eu^{tA} = \eu^{2t}\Bigl(I + tN + \frac{t^2}{2}N^2\Bigr)
= \eu^{2t}\begin{pmatrix}
1 & t & t^2/2\\
0 & 1 & t\\
0 & 0 & 1
\end{pmatrix}.
\]
\end{solution}

\begin{solution}{exo:b2:reduction:8}
$X^k - 1$ 零化 $A$，且在 $\C$ 上完全分裂，有 $k$ 个互异根
$\eu^{2\iu\pi j/k}$：$A$ \omterm{def:b2:reduction:diag}{可对角化}
（\cref{cor:b2:reduction:minpolycrit}），其\omterm{def:b2:reduction:eigen}{特征值}
为 $X^k - 1$ 的根，即 $k$ 次单位根。

若另有 $A$ 相似于对角元全为一的三角矩阵，
所有\omterm{def:b2:reduction:eigen}{特征值}都等于 $1$，$A$ \omterm{def:b2:reduction:diag}{可对角化}，
且唯一\omterm{def:b2:reduction:eigen}{特征值}为 $1$，故 $P\,I\,P^{-1} = I$。
\end{solution}

\begin{solution}{exo:b2:reduction:9}
记 $E = \bigoplus_\lambda E_\lambda(u)$
（由\cref{thm:b2:reduction:diagcrit}）。每个 $E_\lambda(u)$
关于 $v$ 稳定：对 $x \in E_\lambda$，$u(v(x)) = v(u(x)) = \lambda
v(x)$。$v$ 在 $E_\lambda(u)$ 上的限制\omterm{def:b2:reduction:diag}{可对角化}
（见\cref{exo:b2:reduction:4}）：在 $E_\lambda(u)$ 中取一组由
$v$ 的\omterm{def:b2:reduction:eigen}{特征向量}构成的基，对所有 $\lambda$ 合并这些基，
便得 $E$ 的一组基，其向量都是\emph{两者}的\omterm{def:b2:reduction:eigen}{特征向量}：
即既是 $u$ 的（因属于 $E_\lambda(u)$），又是 $v$ 的（由构造）。
\end{solution}

\begin{solution}{exo:b2:reduction:10}
（$\Rightarrow$）设 $u$ \omterm{def:b2:reduction:diag}{可对角化}，$F$ 为稳定子空间。则
$u|_F$ \omterm{def:b2:reduction:diag}{可对角化}（见\cref{exo:b2:reduction:4}）：$F$ 有一组\omterm{def:b2:reduction:eigen}{特征向量}构成的基，
该基在每个全局\omterm{def:b2:reduction:eigen}{特征子空间}内
$E_\lambda$ 中可扩充为 $E_\lambda$ 的一组基（不完全基定理，
在 $E_\lambda$ 内，从 $F$ 的基中落在其中的那部分出发
——注意 $F = \bigoplus_\lambda (F \cap E_\lambda)$，
这由 $u|_F$ \omterm{def:b2:reduction:diag}{可对角化}得出）。新添加的向量张成一个稳定补
（每个向量都落在某个 $E_\lambda$ 中，因而其张成空间
关于 $u$ 稳定）。

（$\Leftarrow$）设 $F = \sum_\lambda E_\lambda(u)$（为稳定
子空间），$G$ 为稳定补。若 $G \neq \{0\}$，则 $\chi_{u|_G}$
在 $\C$ 上完全分裂，故 $u|_G$ 有\omterm{def:b2:reduction:eigen}{特征向量} $x \in G$
（见\cref{thm:b2:reduction:trigonalization}，或直接由根的存在性
得出）；但 $u$ 的每个\omterm{def:b2:reduction:eigen}{特征向量}都属于 $F$，故 $x \in F
\cap G = \{0\}$，矛盾。从而 $G = \{0\}$，$E = F$：
\omterm{def:b2:reduction:eigen}{特征子空间}张满 $E$，即 $u$ \omterm{def:b2:reduction:diag}{可对角化}。
\end{solution}

\begin{solution}{exo:b2:reduction:11}
三角化：$A = PTP^{-1}$，$T = \begin{pmatrix} \lambda & c\\ 0 &
\mu\end{pmatrix}$，$\abs\lambda, \abs\mu < 1$。于是 $A^k =
PT^kP^{-1}$，只需证 $T^k \to 0$。

\emph{互异\omterm{def:b2:reduction:eigen}{特征值}：}由归纳法可得
\[
T^k = \begin{pmatrix}
\lambda^k & c\,\dfrac{\lambda^k - \mu^k}{\lambda - \mu}\\[4pt]
0 & \mu^k
\end{pmatrix},
\]
且每个元素都趋于 $0$（$\abs{\lambda}^k, \abs\mu^k \to 0$）。

\emph{\omterm{def:b2:reduction:eigen}{特征值}相等（$\mu = \lambda$）：}$T = \lambda I + cE_{12}$
且 $T^k = \lambda^k I + k\lambda^{k-1}cE_{12}$，其中元素
$k\lambda^{k-1} \to 0$，因 $\abs\lambda < 1$（几何衰减胜过
多项式增长）。两种情形下均有 $T^k \to 0$（逐元素），从而 $A^k =
PT^kP^{-1} \to 0$（固定的 $P, P^{-1}$ 对矩阵乘法
关于元素连续——乘积的每个元素都是固定的
线性组合）。
\end{solution}

\begin{solution}{exo:b2:reduction:12}
$\ker u$ 的维数为 $n - 1$（秩–零度定理），故 $0$ 是
\omterm{def:b2:reduction:charpoly}{几何重数}为 $n - 1$ 的\omterm{def:b2:reduction:eigen}{特征值}，且 $\chi_u$ 可被
$X^{n-1}$ 整除
（见\cref{def:b2:reduction:charpoly}：\omterm{def:b2:reduction:charpoly}{几何重数} $\leq$ \omterm{def:b2:reduction:charpoly}{代数重数}）。
设 $\chi_u = X^{n-1}(X - \alpha)$，则
$X^{n-1}$ 的系数为 $-\operatorname{tr} u$，故 $\alpha =
\operatorname{tr} u$，于是 $\chi_u = X^{n-1}(X - \operatorname{tr}
u)$。

若 $\operatorname{tr} u \neq 0$，则\omterm{def:b2:reduction:eigen}{特征值}
$\operatorname{tr} u$ 是 $\chi_u$ 的根，故它有一个
\omterm{def:b2:reduction:eigen}{特征向量}；$0$ 与 $\operatorname{tr} u$ 对应的\omterm{def:b2:reduction:eigen}{特征子空间}
维数分别为 $n - 1$ 与 $\geq 1$，其和 $\geq n$，它们
张满 $E$，从而 $u$ \omterm{def:b2:reduction:diag}{可对角化}
（由\cref{thm:b2:reduction:diagcrit}）。若 $\operatorname{tr} u = 0$：
由\cref{exo:b2:linalg:5}，$u^2 = (\operatorname{tr} u)u = 0$
且 $u \neq 0$：$u$ 为非零幂零映射，而\omterm{def:b2:reduction:diag}{可对角化}的幂零映射
为零（见\cref{prop:b2:reduction:nilpotent}），故
不\omterm{def:b2:reduction:diag}{可对角化}。
\end{solution}

\begin{solution}{pb:b2:reduction:1}
\textbf{1.} 前 $k - 1$ 个分量（$Cv_n$ 的分量）为 $u_{n+1},
\dots, u_{n+k-1}$（上对角线的移位），最后一个分量为
$a_0 u_n + \dots + a_{k-1}u_{n+k-1}$。于是 $v_{n+1} = Cv_n$ 对所有
$n$ 成立，当且仅当对所有 $n$ 最后分量相符，即
$(\mathcal R)$ 成立。迭代可得 $v_n = C^nv_0$。

\textbf{2.} 将 $D_k(X) = \det(XI_k - C)$ 按第一列
展开：两个非零元为 $X$（位于 $(1,1)$）与
$-a_0$（位于 $(k,1)$）。第一个余子式形如 $D_{k-1}$，
其系数为 $a_1, \dots, a_{k-1}$；第二个余子式为
对角元为 $-1$ 的上三角矩阵，\omterm{def:b2:linalg:det}{行列式}为 $(-1)^{k-1}$，
符号为 $(-1)^{k+1}$（由位置决定）。对 $k$ 归纳
（基础情形 $k = 1$：$X - a_0$）得
\[
D_k(X) = X\bigl(X^{k-1} - a_{k-1}X^{k-2} - \dots - a_1\bigr) -
a_0 = P(X).
\]
对 $\mu_C$：由 $Q(C^{\mathsf T}) = Q(C)^{\mathsf T}$ 对任意
多项式成立，$C$ 与 $C^{\mathsf T}$ 有相同的\omterm{def:b2:linalg:annihilator}{零化子}，
从而\omterm{def:b2:reduction:polyu}{极小多项式}相同。对 $C^{\mathsf T}$ 而言，
其列满足 $C^{\mathsf T}e_1 = e_2$，\dots，$C^{\mathsf
T}e_{k-1} = e_k$，故 $(e_1, C^{\mathsf T}e_1, \dots, (C^{\mathsf
T})^{k-1}e_1)$ 为标准基，故线性无关。多项式 $Q
\neq 0$ 的次数 $< k$ 时，则有 $Q(C^{\mathsf T})e_1 \neq 0$
（它是基向量的非平凡线性组合），故 $\deg\mu \geq
k$。由 $\mu \mid \chi = P$ 且 $\deg P = k$，得 $\mu_C = P$。

\textbf{3.} 取 $v = (1, \lambda, \dots, \lambda^{k-1})^{\mathsf
T}$：第 $1$ 至 $k-1$ 行（$Cv$ 中）给出 $\lambda, \lambda^2, \dots,
\lambda^{k-1}$，即 $\lambda$ 乘以前 $k - 1$ 个
$v$ 的分量；最后一行给出 $\sum_m a_m\lambda^m = \lambda^k -
P(\lambda) = \lambda^k = \lambda\cdot\lambda^{k-1}$。故 $Cv =
\lambda v$。反之，方程 $(Cx)_i = \lambda x_i$
（$i < k$）即 $x_{i+1} = \lambda x_i$：任一\omterm{def:b2:reduction:eigen}{特征向量}都
与 $v$ 成比例——每个\omterm{def:b2:reduction:eigen}{特征子空间}的维数恰为
$1$。\omterm{def:b2:reduction:diag}{可对角化}当且仅当\omterm{def:b2:reduction:eigen}{特征子空间}维数之和为 $k$
（由\cref{thm:b2:reduction:diagcrit}），即有 $k$ 个互异\omterm{def:b2:reduction:eigen}{特征值}，
亦即 $P$ 有 $k$ 个互异根（\omterm{def:b2:reduction:eigen}{特征值}恰为
$\chi_C = P$ 的根）。

\textbf{4.} 每个 $(\lambda_i^n)_n$ 都满足 $(\mathcal R)$：
$\lambda_i^{n+k} = \lambda_i^n\,\lambda_i^k =
\lambda_i^n\sum_m a_m\lambda_i^m$。无关性：若一个
为零的组合 $\sum_i c_i\lambda_i^n = 0$ 对 $n = 0, \dots, k-1$ 成立，
是关于 $c_i$ 的范德蒙德方程组（见\cref{exo:b2:linalg:11}），
则所有 $c_i = 0$。解空间的维数为 $k$（第 6 题，
其证明初等且与此独立）：$k$ 个自由解
构成一组基，坐标唯一。

\textbf{5.} $P = X^2 - X - 6 = (X - 3)(X + 2)$：通解
$u_n = A\,3^n + B(-2)^n$。初值条件：$A + B = 1$，$3A -
2B = 8$：$A = 2$，$B = -1$：
\[
u_n = 2\cdot 3^n - (-2)^n .
\]
（检验：$u_2 = u_1 + 6u_0 = 14$，且 $2\cdot9 - 4 = 14$。）

\textbf{6.} $P(S)\bigl((u_n)\bigr)$ 为数列 $n \mapsto
u_{n+k} - a_{k-1}u_{n+k-1} - \dots - a_0u_n$：它为零当且仅当
$(\mathcal R)$ 成立，因此解集 $\ker P(S)$ 是一个
子空间。映射 $\ker P(S) \to \C^k$，$u \mapsto (u_0, \dots,
u_{k-1})$ 线性且单射（递推关系确定
$u_{k}, u_{k+1}, \dots$ 由前 $k$ 个值归纳得出），
并且为满射（由任意初始数据递归定义 $u_n$）：
故维数为 $k$。

\textbf{7.} \cref{thm:b2:reduction:kernels} 的证明只用到：
$\C[X]$ 中的贝祖恒等式，以及
同一自同态的多项式彼此可交换这一事实。两者均不涉及
所在空间的维数：引理对
$S \in \mathcal{L}(\mathcal{S})$ 逐字成立，于是
\[
\ker P(S) = \bigoplus_{i=1}^{r}
\ker\,(S - \lambda_i\,\mathrm{id})^{m_i}.
\]

\textbf{8.} 对 $Q \in \C[X]$：$(S -
\lambda)\bigl(Q(n)\lambda^n\bigr)_n$ 的第 $n$ 项为
$Q(n{+}1)\lambda^{n+1} - \lambda Q(n)\lambda^n =
\lambda^{n+1}(\Delta Q)(n)$，其中 $\Delta Q = Q(X{+}1) - Q(X)$
次数为 $\deg Q - 1$（首项相消）。迭代可得
$(S - \lambda)^m\bigl(Q(n)\lambda^n\bigr) =
\bigl(\lambda^{n+m}(\Delta^m Q)(n)\bigr)_n$，且 $\Delta^m Q = 0$
当 $\deg Q \leq m - 1$ 时：右端集合包含于
核中。它是维数为 $m$ 的子空间：序列
$(n^j\lambda^n)_n$，$0 \leq j < m$，线性无关，因为
$\sum_j c_j n^j\lambda^n = 0$ 对所有 $n$ 成立，迫使（两边除以
$\lambda^n \neq 0$）多项式 $\sum_j c_jX^j$ 在
每个 $n \in \N$ 处为零，从而恒为零。反之，$\dim\ker(S -
\lambda)^m \leq m$：展开 $(S - \lambda)^m = \sum_j
\binom mj(-\lambda)^{m-j}S^j$，方程 $(S - \lambda)^m u =
0$ 是 $m$ 阶线性递推式（首项系数
为 $1$），因此 $u$ 由 $u_0, \dots, u_{m-1}$ 确定，方式同
第 6 题。维数相等即得证。

\textbf{9.} 合并第 7 题与第 8 题：每个解都唯一地
分解为 $\ker(S -
\lambda_i)^{m_i}$ 中元素之和，即 $u_n = \sum_i Q_i(n)\lambda_i^n$，其中
$\deg Q_i \leq m_i - 1$；$Q_i$ 唯一，是因为
分解是直和，且在每个求和项中，
$Q_i$ 的系数是关于基
$(n^j\lambda_i^n)_j$ 的坐标（第 8 题）。维数检验：
$\sum_i m_i = k$。

\textbf{10.} $P = X^2 - 4X + 4 = (X - 2)^2$：解为 $(a +
bn)2^n$。初值：$a = 1$，$2(a + b) = 0$，故 $b = -1$：
\[
u_n = (1 - n)\,2^n .
\]
检验：$u_2 = 4u_1 - 4u_0 = -4$，且 $(1 - 2)\cdot4 = -4$。

\textbf{11.} 记 $u_n = \lambda_1^n\bigl(c_1 + \sum_{i\geq2}
c_i(\lambda_i/\lambda_1)^n\bigr)$；每个比值的模均 $< 1$，
故括号趋于 $c_1 \neq 0$：$u_n \sim c_1\lambda_1^n$。
特别地，$u_n \neq 0$（$n$ 充分大时），且
\[
\frac{u_{n+1}}{u_n} =
\lambda_1\,\frac{c_1 + o(1)}{c_1 + o(1)} \longrightarrow
\lambda_1 .
\]

\textbf{12.} 计算：
\[
q(a_{n+1}, b_{n+1}) = (a_n + 2b_n)^2 - 2(a_n + b_n)^2
= -a_n^2 + 2b_n^2 = -q(a_n, b_n).
\]
由 $q(a_0, b_0) = 1 - 2 = -1$，得 $a_n^2 - 2b_n^2 = (-1)^{n+1}$。
从结构上看：$q(a, b) = (a - \sqrt2\,b)(a + \sqrt2\,b)$，且
线性映射 $M$ 将因子 $a + \sqrt2 b$ 乘以 $1 +
\sqrt2$，将因子 $a - \sqrt2 b$ 乘以 $1 - \sqrt2$（计算：
$a_{n+1} + \sqrt2 b_{n+1} = (1 + \sqrt2)(a_n + \sqrt2 b_n)$）；
乘积每一步乘以 $(1 + \sqrt2)(1 - \sqrt2) = -1 =
\det M$。

\textbf{13.} 由于 $a_n^2 - 2b_n^2 = (a_n - \sqrt2 b_n)(a_n +
\sqrt2 b_n) = (-1)^{n+1}$，
\[
\Bigl|\frac{a_n}{b_n} - \sqrt2\Bigr|
= \frac{\abs{a_n^2 - 2b_n^2}}{b_n(a_n + \sqrt2 b_n)}
= \frac{1}{b_n(a_n + \sqrt2 b_n)} \leq \frac1{2b_n^2},
\]
利用 $a_n \geq b_n \geq 1$（归纳法：两者都递增），故 $a_n +
\sqrt2 b_n \geq (1 + \sqrt2)b_n \geq 2b_n$。$M$ 的\omterm{def:b2:reduction:eigen}{特征值}：
$\chi_M = X^2 - 2X - 1$，根为 $1 \pm \sqrt2$；由于 $(a_0, b_0)$
在主\omterm{def:b2:reduction:eigen}{特征向量}上有非零分量（各分量均
为正），$b_n \sim c(1 + \sqrt2)^n$，其中 $c > 0$（第 11 题）。
于是误差 $\asymp (1 + \sqrt2)^{-2n} = (3 +
2\sqrt2)^{-n}$：几何衰减，公比为 $1/(3 + 2\sqrt2) = 3
- 2\sqrt2 \approx 0.172$。

\textbf{14.} (a) 由第 9 题：$\abs{u_n} \leq \sum_i
\abs{Q_i(n)}\abs{\lambda_i}^n \leq \bigl(\sum_i
\abs{Q_i(n)}\bigr)\rho^n$，且每个 $\abs{Q_i(n)} \leq C_i
n^{m_i - 1} \leq C_i n^{m-1}$，对 $n \geq 1$ 成立，将常数相加即可。
(b) 令 $\rho' = \max_{i \geq 2}\abs{\lambda_i} <
\abs{\lambda_1}$，$d = \deg Q_1$，首项系数为 $c
\neq 0$。于是 $u_n = Q_1(n)\lambda_1^n + R_n$，其中 $\abs{R_n}
\leq Cn^{m-1}\rho'^n$，且
\[
\frac{R_n}{Q_1(n)\lambda_1^n} = O\Bigl(n^{m-1-d}
\bigl(\rho'/\abs{\lambda_1}\bigr)^n\Bigr) \longrightarrow 0
\]
（几何衰减压倒多项式增长）。因此
\[
u_n \sim Q_1(n)\,\lambda_1^n \sim c\,n^d\lambda_1^n,
\qquad
\frac{u_{n+1}}{u_n} \longrightarrow \lambda_1
\quad\text{（因为 } Q_1(n{+}1)/Q_1(n) \to 1\text{）}.
\]
检验第 10 题：对 $u_n = (1-n)2^n$，比值为
\[
\frac{(-n)2^{n+1}}{(1-n)2^n} = 2\,\frac{-n}{1-n}
\longrightarrow 2 = \lambda_1 .
\]

\textbf{15.} 对 $n$ 用归纳法。当 $n = 1$ 时，$A_{ij}$ 计数
长度为 $1$ 的道路。归纳步：长度为 $n + 1$ 的从 $i$ 到
$j$ 的道路，可分解为长度为 $n$ 的、从 $i$ 到某顶点 $\ell$ 的道路，
再接一条边 $\ell j$：
\[
\#\{\text{道路}\} = \sum_{\ell} (A^n)_{i\ell}A_{\ell j} =
(A^{n+1})_{ij}.
\]

\textbf{16.} $J = 3\Pi$，其中 $\Pi = J/3$ 是
$\operatorname{Vect}(1,1,1)$ 上、沿平面 $x + y + z = 0$ 的投影
（$\Pi^2 = \Pi$，因为 $J^2 = 3J$）。于是 $A = J - I = 2\Pi - (I -
\Pi)$，由于 $\Pi$ 与 $I - \Pi$ 是互补的
投影，
\[
A^n = 2^n\,\Pi + (-1)^n (I - \Pi),
\qquad\text{即}\qquad
(A^n)_{ij} = \frac{2^n}3 + (-1)^n\Bigl(\delta_{ij} -
\frac13\Bigr),
\]
由此得到题中两个公式。当 $n = 2$ 时：对角线元素为
$(4 + 2)/3 = 2$（道路 $i \to \ell \to i$，对两个邻点 $\ell$）；
非对角线元素为 $(4 - 1)/3 = 1$（唯一道路 $i \to
\ell \to j$ 经过第三个顶点）。

\textbf{17.} 设 $w_n^{(0)}, w_n^{(1)}$ 分别为长度为 $n$ 的允许词中
以 $0$、$1$ 结尾的个数。添加一个字母：$0$
可以接在任何字母之后，$1$ 只能接在 $0$ 之后：
\[
\begin{pmatrix} w_{n+1}^{(0)}\\ w_{n+1}^{(1)}\end{pmatrix}
= \begin{pmatrix} 1 & 1\\ 1 & 0\end{pmatrix}
\begin{pmatrix} w_n^{(0)}\\ w_n^{(1)}\end{pmatrix}.
\]
求和得 $w_{n+2} = w_{n+1} + w_n$（或：按首字母
分类讨论）。由 $w_1 = 2$，$w_2 = 3$，归纳得 $w_n = F_{n+2}$
（$F_3 = 2$，$F_4 = 3$，递推式相同）。增长率：$X^2 - X - 1$ 的根为
$\varphi > \abs\psi$，
（见\cref{exo:b2:reduction:5}），而 $\varphi$ 的分量
非零（$w_n$ 为正且 $\psi^n \to 0$），故由第 11 题
得 $w_{n+1}/w_n \to \varphi = \frac{1 + \sqrt5}2$。

\textbf{18.} $A = \begin{psmallmatrix} 0&1&0\\ 1&0&1\\
0&1&0\end{psmallmatrix}$。验证：
\[
A(1, \pm\sqrt2, 1)^{\mathsf T}
= (\pm\sqrt2, 2, \pm\sqrt2)^{\mathsf T}
= \pm\sqrt2\,(1, \pm\sqrt2, 1)^{\mathsf T},
\qquad
A(1, 0, -1)^{\mathsf T} = 0 :
\]
\omterm{def:b2:reduction:eigen}{特征值}为 $\sqrt2, -\sqrt2, 0$（$= 2\cos\frac\pi4,
2\cos\frac{3\pi}4, 2\cos\frac\pi2$）。将 $e_1$ 在特征基
上分解并读取第三个坐标，或利用对称性，取
$v_\pm = (1, \pm\sqrt2, 1)$，$v_0 = (1, 0, -1)$，可验证 $e_1
= \frac14 v_+ + \frac14 v_- + \frac12 v_0$，因此对 $n \geq 1$ 有
\[
(A^n)_{13} = \Bigl(\tfrac14(\sqrt2)^n v_+ +
\tfrac14(-\sqrt2)^n v_- + 0\Bigr)_{\!3}
= \frac{(\sqrt2)^n + (-\sqrt2)^n}{4},
\]
当 $n$ 为奇数时为零（二部图：两端点之间的距离为偶数），
且 $2\cdot 2^{n/2}/4 = 2^{n/2 - 1}$，对偶数 $n$ 成立。当 $n = 4$ 时：
$2^{1} = 2$，与两条道路 $1\,2\,1\,2\,3$ 和
$1\,2\,3\,2\,3$ 相符。

\textbf{19.} 长度为 $n$ 的、从 $i$ 出发的闭道路数为
$(A^n)_{ii}$；对 $i$ 求和得 $\operatorname{tr}(A^n)$。
将 $A$ 上三角化（在 $\C$ 上），$A^n$ 为上三角矩阵，对角元为 $\lambda_i^n$：
$\operatorname{tr}(A^n) = \sum_i\lambda_i^n$。
三角形：$\operatorname{tr}(A^n) = 3\,\frac{2^n + 2(-1)^n}3 =
2^n + 2(-1)^n = 2^n + (-1)^n + (-1)^n$：谱为 $\{2, -1,
-1\}$，与第 16 题一致。

\textbf{20.} $W^{\mathsf T}$ 的列为：$W^{\mathsf T}e_i =
e_{i-1}$（$i \geq 1$），且 $W^{\mathsf T}e_0 = e_{k-1}$；
按 $e_0, e_1, \dots$ 的顺序重新编号，这恰好是
$X^k - 1$ 的友矩阵（$a_0 = 1$，其余 $a_m = 0$）。
第 2 题：$\chi_{W} = \chi_{W^{\mathsf T}} = X^k - 1 =
\mu_{W}$。根 $\omega^j$（$j = 0, \dots, k-1$）是 $k$
个互异的 $k$ 次单位根：$W$ \omterm{def:b2:reduction:diag}{可对角化}（第 3 题
，或见\cref{exo:b2:reduction:8}：$W^k = I$）。\omterm{def:b2:reduction:eigen}{特征向量}：$Wf_j
= \sum_m \omega^{-jm}e_{m+1} = \sum_{m'}\omega^{-j(m'-1)}e_{m'}
= \omega^j f_j$.

\textbf{21.} 循环矩阵是 $W$ 的多项式，而
同一固定矩阵的多项式彼此可交换。每个 $f_j$ 都是
所有幂的\omterm{def:b2:reduction:eigen}{特征向量}：$W^m f_j = \omega^{jm}f_j$，故
\[
Cf_j = \sum_m c_m\omega^{jm} f_j = \widehat c(\omega^j)\,f_j :
\]
基 $(f_0, \dots, f_{k-1})$（线性无关：由互异的
$\omega^{-j}$ 构成的范德蒙矩阵，见\cref{exo:b2:linalg:11}）同时对角化
每个循环矩阵，\omterm{def:b2:reduction:eigen}{特征值}如所述。

\textbf{22.} \omterm{def:b2:linalg:det}{行列式}等于\omterm{def:b2:reduction:eigen}{特征值}的乘积
（对角化）：$\det C = \prod_{j}\widehat c(\omega^j)$。当 $k
= 3$，$c_0 = a$，$c_1 = b$，$c_2 = c$，且 $\omega = j =
\eu^{2\iu\pi/3}$:
\[
\det C = (a + b + c)(a + bj + cj^2)(a + bj^2 + cj^4),
\]
又 $j^4 = j$：这恰好就是
\cref{exo:b2:linalg:8} 中的因式分解。

\textbf{23.} $M = \frac12(W + W^{-1})$ 是循环矩阵（$W^{-1} =
W^{k-1}$），\omterm{def:b2:reduction:eigen}{特征值}为 $\frac12(\omega^j + \omega^{-j}) =
\cos\frac{2\pi j}k$，\omterm{def:b2:reduction:eigen}{特征向量}仍取同一组基 $f_j$。坐标：记
$x^{(0)} = \sum_j \alpha_j f_j$。$f_j$ 的坐标之和为
$\sum_m \omega^{-jm}$，其值为 $k$（仅当 $j = 0$
时），其余情形为 $0$（公比 $\omega^{-j} \neq 1$ 的等比求和）。
对 $x^{(0)}$ 的坐标求和：$\sum_m x^{(0)}_m =
\alpha_0\,k$，故 $\alpha_0 = \frac1k\sum_m x^{(0)}_m$，即均值。

\textbf{24.} $x^{(n)} = M^nx^{(0)} = \sum_j
\alpha_j\cos^n\bigl(\tfrac{2\pi j}k\bigr)f_j$。当 $k$ 为奇数时，
$\abs{\cos(2\pi j/k)} < 1$ 对每个 $j \neq 0$ 成立（角度
永不为 $0$ 或 $\pi$），故除 $j = 0$ 外各项趋于 $0$：
$x^{(n)} \to \alpha_0 f_0$，即常值向量，其各分量都等于
均值——奇数环上取平均会使各值趋于一致。当 $k = 4$ 时，
\omterm{def:b2:reduction:eigen}{特征值}为 $1, 0, -1, 0$：$j = 2$ 对应的项
$\alpha_2(-1)^nf_2$，其中 $f_2 = (1, -1, 1, -1)^{\mathsf T}$
永远振荡。障碍在于交错均值：
将 $x^{(0)}$ 的各坐标乘以 $(-1)^m$ 并
求和，同样用等比求和计算得 $\sum_m
(-1)^mx^{(0)}_m = 4\alpha_2$：该过程收敛当且仅当
$x^{(0)}_0 - x^{(0)}_1 + x^{(0)}_2 - x^{(0)}_3 = 0$，此时
收敛于均值。

\textbf{25.} 友矩阵将一个
$k$ 阶标量递推式化为一阶向量递推式，使得
闭式公式成为关于 $C^n$ 的结论——约化理论的
用武之地（第 1--5 题）。核分解引理纯属
多项式代数（贝祖恒等式加交换性），故能分解
$\ker P(S)$，尽管 $\mathcal{S}$ 是无限维的
（第 7--9 题）。主\omterm{def:b2:reduction:eigen}{特征值}决定增长，因为
其余各项按几何速度变得可忽略，
在归一化之后——这也正是佩尔误差按
主根的平方速率衰减的原因（第 11--14 题）。
邻接矩阵的幂计数游走，因为矩阵乘法
对中间顶点求和，因此谱计数闭游走
（第 15--19 题）。可交换矩阵有共同的特征基，而
于是一组傅里叶基便能对角化整个循环矩阵代数，
一举完成（第 20--24 题）。要点：线性递推的
基本定理（第 9 题）；对于非负
矩阵，主根（如 $\varphi$ 或 $1 +
\sqrt2$）自动为实数、为正且为单根，其原因在于
佩龙--弗罗贝尼乌斯定理，在本科三年级分册中证明。
\end{solution}
